\section{Related Work}
\label{sec:related}

\para{Semantic and confidentiality leakage.} \citet{gonen2025semantic} define semantic leakage as excess influence from an unrelated prompt concept and demonstrate it with matched controls across languages and generation tasks. \citet{holtzman2026secret} adapt this idea to explicit secrecy: models avoid literal mentions but write stories that disclose the hidden word through thematic choices. Their active-hiding and decoy conditions show that avoidance can invert or redirect the signal. We build directly on their 15-word, order-balanced story protocol, but ask whether a secret-blind outline changes leakage beyond the effect of extra context. Privacy benchmarks such as \textsc{ConfAIde} measure direct, norm-dependent disclosure and emphasize the need to evaluate privacy jointly with utility \citep{mireshghallah2023confaide}. More recently, \citet{fairoze2026inadvertent} recover protected context from benign outputs with adaptive decoders and connect stronger suppression to larger leakage. \citet{mann2025whitebear} likewise find forbidden-token rebound under cognitive load. Unlike these studies, our primary endpoint concerns reader-visible influence in long-form stories rather than direct disclosure, adversarially trained decoding, or next-token probability.

\para{Planning and narrative generation.} Residual-state probes can predict future response length, answers, and other attributes before generation begins \citep{dong2025emergent}. Myopic training experiments distinguish deliberate preparation for future tokens from semantic ``breadcrumbs'' that are useful both now and later \citep{wu2024plan}. This distinction motivates our central hypothesis: an outline may reduce repeated preparation without removing ordinary semantic association. Narrative work shows that iterative plot construction followed by prose elaboration improves suspense and consistency \citep{xie2024suspenseful}, but does not isolate what happens when the prose model holds information absent from the plan. We therefore compare a secret-blind outline with both no outline and length-matched irrelevant text. Private working memory addresses a related but different problem---maintaining a hidden state across public interaction \citep{baldelli2026hangman}. Our outline is public conditioning text, not a private memory channel.

\para{Internal secret representations.} Model-organism studies recover hidden concepts through behavioral elicitation, logit-lens readouts, and sparse-autoencoder features \citep{cywinski2025latent,cywinski2025eliciting}. These results establish accessibility, not that a recovered direction causes natural leakage. Behavioral branch probes also reveal that apparently stable internal targets can drift during interaction \citep{luo2026beliefs}. We therefore learn directions only from independent carrier prompts, inspect their trajectory on held-out story prose, and require a causal intervention to preserve basic generation utility. This utility gate is central: an intervention that removes most of the story cannot distinguish selective information removal from generic model damage.

